% Chapter 7: Discussion

\chapter{Discussion}
\label{ch:discussion}

Chapter~\ref{ch:results} reported the results. This chapter interprets RQ1 and
RQ2 in turn, brings the two stages together, relates the findings to the three
strands reviewed in Chapter~\ref{ch:related-work}, and then states what the
evidence means for faithfulness, its limitations, and future work.

\section{Interpreting RQ1: answer behaviour under changes to visual evidence}
\label{sec:discussion-rq1}
\label{sec:discussion-descriptions}

Accuracy under mismatch (0.435--0.519) was lower than accuracy under
grey (0.483--0.626) in every arm. Grey takes the scene away and leaves the
question, the choices, and any description intact. Mismatch leaves a full,
coherent scene in place, but it belongs to a different record. If the model were
simply falling back on textual priors once the image stopped being useful, the
two conditions should not separate in this direction: an unrelated image would
be no worse than no image at all. That the wrong image is worse than no image is
consistent with the model treating whatever it is shown as evidence about the
question, and being misled when that evidence is irrelevant.

The two conditions also differ in input format. Mismatch supplies a full-size
photograph of the kind the model was adapted on, while grey supplies a small
uniform square that matches no frame in the dataset
(Section~\ref{sec:setup-conditions}). An input that is easy to disregard and one
that demands to be read are not the same
manipulation, so this design cannot separate misleading evidence from a
difference in input format. Both explanations concern observable behaviour.

Mirror left aggregate Stage~1 accuracy close to source, but 5.24--7.50\% of
records moved between correct and wrong states. A stable aggregate accuracy can
therefore conceal substantial movement for individual records.

Mirror preserves scene content while changing its orientation, so it provides a
comparison with less visual disruption than removal, corruption, or replacement.
Horizontal reversal can
alter text placement, left-right relations, and pixel organisation.
The observed churn makes that limitation visible. The analysis does not
identify which question types account for the movement.

The description arms lost less Stage~1 accuracy under the interventions, and
under grey the difference sat mainly in fewer losses from correct to wrong rather than
more gains from wrong to correct.
The polygon results differed across comparisons. Differences between the polygon
arms and their plain counterparts were small at source, and the \texttt{point}
against \texttt{plain} difference changed direction across the Stage~1
interventions.

Stage~1 answer behaviour therefore responds to the content of the image as well
as to its presence, and a fixed description partly buffers that response.

\section[Interpreting RQ2: rationale sensitivity and answer conditioning]%
  {Interpreting RQ2: rationale sensitivity and \mbox{answer} conditioning}
\label{sec:discussion-rq2}

Under gold conditioning the rationale text moved for every intervention.
Replacement by an unrelated scene moved it most, removal next, corruption and
added noise less, and reversal least. Arms carrying a description moved less when grey took the
scene away. \texttt{point\_desc} had higher mean Drift than
\texttt{plain\_desc} under all five interventions
(Table~\ref{tab:rq2-drift}), but that comparison was not
planned, so no interval is reported and no causal reading follows: under
mismatch the arms do not receive the same kind of image
(Section~\ref{sec:setup-conditions}).

Gold conditioning held the supplied answer fixed, which made the image
intervention interpretable. The 200-record wrong-answer control then changed
that answer while retaining the source image. Drift from the wrong answer
exceeded matched grey Drift in all four arms, and the direction survived both
the reasoning span and final rationales with answer choices masked. Rationale
generation therefore responded to the answer it was asked to justify as well as
to the image. That comparison is specific to matched grey removal; the other
conditions and anchors have different sample scopes and support no inferential
ranking against answer change.

The predicted-answer regime adds a further observation. When Stage~2
wrote a rationale for an answer Stage~1 had selected, a judge receiving only
text could often recover that answer, right or wrong. In \texttt{plain\_desc}
the conditional rate was almost the same for correct and wrong predictions,
0.5378 against 0.5331 (Appendix~\ref{sec:app-pred-a}). Coherence with a supplied
answer therefore does not establish that the answer is correct. Coverage of usable
rationales ranged from 39.84\% to 97.02\%, so a rate computed over successful
generations alone would hide the records lost. The regime needs the
conditional and the pipeline view together, and these readings stay inside its
own scope (Section~\ref{sec:setup-pred-a}).

Stage~2 rationale text therefore has two observable dependencies. It moves when
the image changes, and when the supplied answer changes.

\section{Relationship between Stage~1 and Stage~2}
\label{sec:discussion-stages}

The two stages showed the same broad intervention ordering even though their
outputs and model heads differed. The comparison is ordinal. An accuracy gap is
a change in the share of correct answers, while Drift is a change in text
similarity among usable rationale pairs. Their magnitudes cannot be placed on
one scale. The common ordering is therefore evidence of agreement between the
stages across interventions. It does
not establish a shared effect size.

Descriptions reduced disruption at both stages. Their buffering effect is
consistent with an alternative information route through text.
Explicit polygon pointing did not yield a consistent independent advantage.
The most consistent robustness pattern was associated with descriptions.

Because each arm was trained once with its own inputs, the comparison describes four
separately trained systems (Section~\ref{sec:problem-scope}). The experiment does
not identify the mechanism or establish stronger visual grounding.

\section{How the findings address gaps in prior work}
\label{sec:discussion-priorwork}

Chapter~\ref{ch:related-work} reviewed visual question answering, rationale
generation, and behavioural evaluation of explanations. The findings address a
specific gap in each strand, within one model family, one VCR population, and
one adaptation per arm.

\subsection{Visual Question Answering}
\label{sec:discussion-priorwork-vqa}

This thesis addresses the visual question answering gap by applying five
controlled image interventions to the same VCR records and evaluating both
answer and rationale behaviour. Descriptions reduced disruption at both stages,
while polygon overlays produced no consistent independent advantage. These
findings extend work showing shortcuts across questions and detected image
content
\citep{DBLP:conf/iccv/DancetteCTC21,DBLP:conf/emnlp/SiMZLL0CWZ22},
that verbal descriptions can substitute for or compete with visual input, and
that a region label does not guarantee use of the linked image area
\citep{HakimovS23,DBLP:conf/cvpr/DengCCH25,LiGWCNK24}.

\subsection{Rationale Generation}
\label{sec:discussion-priorwork-rationales}

This thesis addresses the rationale generation gap by testing answer
conditioning directly. In the wrong-answer control, changing the supplied answer
while keeping the source image fixed produced greater Drift than matched grey
replacement in every arm. PRED-A also showed that coherence with the supplied
answer can occur when that answer is wrong. These findings extend work on
answer-conditioned rationale systems
\citep{DBLP:conf/cvpr/ParkHARSDR18,DBLP:conf/emnlp/MarasovicBPLSC20}
and on written reasoning that follows a prompt cue without mentioning it
\citep{DBLP:conf/nips/TurpinMPB23}. The control covers 200 records, and PRED-A
coverage varies across arms.

\subsection{Behavioural Evaluation of Explanations}
\label{sec:discussion-priorwork-behavioural}

This thesis addresses the behavioural evaluation gap by extending intervention
logic to visual evidence in free-form rationale generation. It keeps the record
and supplied gold answer fixed, changes the image condition, and measures paired
rationale change with BERTScore Drift. Applying the same image conditions at
both stages supports an ordinal cross-stage synthesis. Prior work intervened on
input spans, attention distributions, or written chains of reasoning
\citep{DeYoungEtAl20ERASER,JainWallace19Attention,DBLP:journals/corr/abs-2307-13702}.
Related multimodal work compares image and text contributions to answers and
explanations \citep{ParcalabescuFrank25VLM}. The specific contribution here is a
matched VCR setup using the same records and one image-intervention battery at
both stages, holding the gold answer fixed in primary Stage~2 and interpreting
rationale change with empirical Drift references.

\section{What the findings mean for faithfulness}
\label{sec:discussion-faithfulness}

% Without this the page ends on the Section 7.6 heading and its one-line
% introduction, leaving the seven bullets to open the next page on their own.
% Shortening this page carries the heading and its introduction across with them.
\enlargethispage{-4\baselineskip}
Stage~1 answers and Stage~2 rationale text respond systematically to changes in
visual evidence, and the rationale text also responds to the answer it is given.
These are behavioural signals relevant to whether a displayed rationale tracks
the available evidence.

The measures do not observe the causal computation that produced a particular
answer. BERTScore Drift records textual change and has no universal faithfulness
threshold. A low value can arise from stable content, shared phrasing, or other
properties of the generator. A high value records greater text change but does not identify why. The study
therefore supports comparative claims about
behavioural sensitivity. It does not support labels for individual rationales.

The reference using rationales from different records illustrates the same
point. It locates the observed values on this dataset's raw BERTScore scale. It
is an empirical reference produced by one fixed pairing.

\section{Limitations}
\label{sec:discussion-limitations}

The main limitations are:

\begin{itemize}
  \item \textbf{Scope.} The study uses one model family, one VCR population,
        one adaptation seed per arm, and separate checkpoints for each arm.

  \item \textbf{Evaluation.} The reported values are performance on the
        validation set rather than estimates from a separate untouched test
        set: the same records supported checkpoint assessment and reporting.
        Stage~1 selection and Stage~2 BLEU monitoring used outputs from the
        source condition only. The bootstrap intervals cover variation across
        records and source frames, excluding uncertainty from checkpoint
        choice, training seed, and adaptation.

  \item \textbf{Descriptions.} The GPT-4o API alias was not pinned to a dated
        model snapshot. The descriptions were generated from polygon images,
        were not independently checked for factual accuracy or task coverage,
        and were truncated to 256 characters. A description error could cause
        or mask apparent textual substitution. In the two description arms,
        the Stage~2 caption block differs between training and primary
        generation (Appendix~\ref{sec:app-prompts}), creating an additional
        train/inference mismatch specific to those arms. This additionally
        confounds the cross-arm Stage~2 description comparison. Within-arm
        source-versus-intervention pairs remain internally matched because
        both members use the same generation prompt.

  \item \textbf{Training loss.} Stage~2 training loss included prompt tokens
        (Section~\ref{sec:setup-prompts}), a training-design limitation whose
        effect on formatting or echo behaviour was not separately evaluated.

  \item \textbf{Visible artefacts.} Some frames contain distributor branding or
        subtitles. Mask, noise, and mirror preserve or partly preserve these
        cues, while grey and mismatch remove them. The interventions therefore
        do not isolate scene content from every other visible cue.

  \item \textbf{Stage comparison and deployment.} Stage~1 and Stage~2 use
        different outcomes and hardware environments, so only qualitative
        agreement between them is supported. Primary Stage~2 uses the gold
        answer. A deployed pipeline would justify the answer selected by
        Stage~1.

  \item \textbf{Output format.} Tolerant parsing recovered usable text from
        some generations that did not follow the requested format. Strict
        formatting failures were concentrated in \texttt{point\_desc}.

  \item \textbf{Transition interpretation.} W$\rightarrow$C transitions can
        reflect choice artefacts or ambiguity and therefore do not establish
        corrected reasoning.
\end{itemize}

\section{Future work}
\label{sec:discussion-future}

Future work should:

\begin{itemize}
  \item replicate the intervention battery across model families, datasets,
        adaptation seeds, and prompt families;

  \item use independent human judgements to relate Drift values to perceived
        changes in rationale content;

  \item vary description availability during training and evaluation to
        separate the buffering effect from the trained arm; and

  \item use mechanistic or causal interventions to test the internal evidence
        routes suggested by the behavioural findings.
\end{itemize}
